% ============================================================
\chapter{Applications: Finance, Economics, Meteorology}
\label{ch:applications}
% ============================================================

\section{Motivating Example}
The tools developed in this course apply directly to concrete problems: portfolio management, macroeconomic forecasting, and climate modeling.

\section{Finance: Volatility Modeling}\index{finance}

\subsection{Value at Risk (VaR)}\index{VaR}

\begin{definition}[VaR]
The \textbf{Value at Risk} at level $\alpha$ is the $\alpha$-quantile of the loss distribution:
\[
P(L_t > \text{VaR}_\alpha) = \alpha.
\]
With a GARCH model: $\text{VaR}_\alpha = -\hat\mu_t + \hat\sigma_t\,q_\alpha$, where $q_\alpha$ is the quantile of the innovation distribution.
\end{definition}

\subsection{Application}

\begin{codeblock}[title=\textbf{Python --- VaR with GARCH}]
\begin{minted}{python}
import numpy as np
import pandas as pd
from arch import arch_model

# Simuler des rendements type financier
np.random.seed(42)
n = 2000
omega, alpha, beta = 0.01, 0.08, 0.90
sigma2 = np.zeros(n)
r = np.zeros(n)
sigma2[0] = omega / (1 - alpha - beta)
for t in range(1, n):
    sigma2[t] = omega + alpha * r[t-1]**2 + beta * sigma2[t-1]
    r[t] = np.sqrt(sigma2[t]) * np.random.standard_t(5)

# Estimer GARCH(1,1) avec innovations t de Student
am = arch_model(r * 100, vol='Garch', p=1, q=1, dist='t')
res = am.fit(disp='off')

# VaR 1%
forecasts = res.forecast(horizon=1)
sigma_next = np.sqrt(forecasts.variance.values[-1, 0])
from scipy.stats import t as t_dist
nu = res.params['nu']
VaR_01 = sigma_next * t_dist.ppf(0.01, nu)
print(f"VaR 1% = {VaR_01:.4f}%")
\end{minted}
\end{codeblock}

\section{Economics: Macroeconomic Forecasting}\index{economics}

\subsection{GDP and Leading Indicators}

\begin{method}
\begin{enumerate}
\item Load quarterly series (GDP, consumption, investment).
\item Test for stationarity (ADF). Difference if necessary.
\item Estimate a VAR or a seasonal ARIMA.
\item Evaluate by out-of-sample validation.
\end{enumerate}
\end{method}

\begin{codeblock}[title=\textbf{Python --- Macro Forecasting}]
\begin{minted}{python}
import pandas as pd
import numpy as np
from statsmodels.tsa.api import VAR
from statsmodels.tsa.stattools import adfuller

# Exemple avec données simulées (remplacer par FRED/Eurostat)
np.random.seed(0)
n = 200
gdp = np.cumsum(np.random.normal(0.5, 1, n))      # trend + bruit
cpi = 0.3 * gdp + np.cumsum(np.random.normal(0, 0.5, n))

df = pd.DataFrame({'GDP_growth': np.diff(gdp),
                    'Inflation': np.diff(cpi)})

# VAR
model = VAR(df)
results = model.fit(ic='bic')
print(f"Ordre sélectionné: VAR({results.k_ar})")
print(results.summary())

# Prévision
forecast = results.forecast(df.values[-results.k_ar:], steps=8)
print("Prévision PIB (8 trimestres):", forecast[:, 0])
\end{minted}
\end{codeblock}

\section{Meteorology: Temperatures and Seasonality}\index{meteorology}

\begin{codeblock}[title=\textbf{Python --- Temperature Series}]
\begin{minted}{python}
import numpy as np
import pandas as pd
import matplotlib.pyplot as plt
from statsmodels.tsa.seasonal import seasonal_decompose
from statsmodels.tsa.statespace.sarimax import SARIMAX

# Simuler des températures mensuelles
np.random.seed(42)
n_years = 30
t = np.arange(n_years * 12)
seasonal = 10 * np.sin(2 * np.pi * t / 12)
trend = 0.01 * t  # réchauffement
noise = np.random.normal(0, 2, len(t))
temp = 15 + trend + seasonal + noise

df = pd.DataFrame({'temp': temp},
    index=pd.date_range('1990-01', periods=len(t), freq='ME'))

# Décomposition
decomp = seasonal_decompose(df['temp'], model='additive', period=12)
decomp.plot()
plt.savefig('ch11_decomp.pdf')
plt.show()

# SARIMA
train = df.iloc[:-24]
test = df.iloc[-24:]
model = SARIMAX(train, order=(1,0,1),
                seasonal_order=(1,1,1,12))
results = model.fit(disp=False)
pred = results.forecast(steps=24)

rmse = np.sqrt(np.mean((test['temp'] - pred.values)**2))
print(f"RMSE prévision 24 mois: {rmse:.2f}°C")
\end{minted}
\end{codeblock}

\section{Methodological Summary}

\begin{center}
\begin{tabular}{lll}
\toprule
\textbf{Domain} & \textbf{Typical Model} & \textbf{Main Challenge} \\
\midrule
Finance & GARCH, EGARCH & Volatility, heavy tails \\
Macroeconomics & VAR, VECM, ARIMA & Causality, cointegration \\
Meteorology & SARIMA, state space & Seasonality, trend \\
Epidemiology & SIR + state space & Nonlinearity \\
Signal processing & Spectral analysis & Hidden periodicities \\
\bottomrule
\end{tabular}
\end{center}

\section{Exercises}

\begin{exercise}[$\star$]
Compute the 5\% VaR for a GARCH(1,1) with $\hat\sigma_{t+1}=2\%$ and Gaussian innovations.
\end{exercise}

\begin{exercise}[$\star\star$]
Compare SARIMA and Holt--Winters exponential smoothing forecasts on monthly temperatures. Which is more accurate?
\end{exercise}

\begin{exercise}[$\star\star\star$ -- Project]
Choose a real-world dataset in one of the three domains. Apply the complete methodology: exploration, modeling (at least 3 models), diagnostics, out-of-sample forecasting, comparison. Write a 10-page report.
\end{exercise}

\begin{keyformulas}
\begin{align*}
\text{VaR}_\alpha &= -\hat\mu_t + \hat\sigma_t\,q_\alpha \\
\text{RMSE} &= \sqrt{\frac{1}{H}\sum(X_{n+h}-\hat{X}_{n+h})^2} \\
\text{DM} &= \frac{\bar{d}}{\sqrt{\hat\sigma_d^2/H}}\dto\mathcal{N}(0,1)
\end{align*}
\end{keyformulas}
